\documentclass[a4paper,12pt]{article}

% Pacchetti essenziali
\usepackage[utf8]{inputenc}
\usepackage[T1]{fontenc}
\usepackage[italian]{babel}
\usepackage{geometry}
\geometry{top=2.5cm, bottom=2.5cm, left=2.5cm, right=2.5cm}
\usepackage{hyperref}
\hypersetup{
	colorlinks=true,
	linkcolor=blue,
	urlcolor=blue,
	pdftitle={Manuale Utente - Portale Cori San Vito},
	pdfsubject={Manuale utente}
}
\usepackage{graphicx} % opzionale per eventuali screenshot
\usepackage{enumitem}
\usepackage{booktabs}
\usepackage{array}
\usepackage{xcolor}

% Comandi personalizzati
\newcommand{\roleadmin}{\texttt{admin}}
\newcommand{\roledirettore}{\texttt{direttore}}
\newcommand{\roleresponsabile}{\texttt{responsabile}}
\newcommand{\rolecorista}{\texttt{corista}}
\newcommand{\rolestrumentista}{\texttt{strumentista}}

\begin{document}
	
	\begin{titlepage}
		\centering
		{\Huge\bfseries Manuale Utente\par}
		\vspace{0.5cm}
		{\Large Portale dei Cori San Vito\par}
		\vspace{1cm}
		{\large Versione 1.0 -- \today\par}
		\vspace{1.5cm}
		\includegraphics[width=0.3\textwidth]{logo.png} % se disponibile
		\vfill
		{\large Parrocchia di San Vito e C.M. -- Spinea (VE)\par}
		\vspace{0.3cm}
		{\large \url{https://github.com/corisanvito/corisanvito.github.io}\par}
	\end{titlepage}
	
	\tableofcontents
	\newpage
	
	\section{Introduzione}
	
	Il portale \textbf{Cori San Vito} è il sito web ufficiale dei cori della parrocchia di San Vito e C.M. a Spinea (VE).  
	Lo scopo del progetto è quello di riunire tre cori (\textit{Coro delle 10}, \textit{Coro delle 11:15} e \textit{Coro Estivo}) in un unico portale, mettendo a disposizione di tutti i coristi e dei responsabili gli strumenti necessari per la gestione quotidiana del gruppo.
	
	Il portale è composto da due aree principali:
	
	\begin{itemize}
		\item \textbf{Area pubblica} (accessibile a tutti): home page di selezione del coro, repertorio canti con ricerca, pagina di richiesta nuovi canti, pagine di dettaglio per ogni coro.
		\item \textbf{Area riservata} (\texttt{/portale}): accessibile solo dopo il login, con funzionalità differenziate in base al ruolo.
	\end{itemize}
	
	\section{Come si accede al portale}
	
	L'accesso all'area riservata avviene tramite la pagina \texttt{/portale/login.html}.  
	Per effettuare il login è necessario disporre di un account e del token di autenticazione fornito dai responsabili.
	
	\begin{important}
		Il token viene gestito automaticamente dal sistema (\texttt{script/auth-guard.js}) e viene verificato ad ogni caricamento di una pagina riservata.
	\end{important}
	
	\section{Ruoli e permessi}
	
	Il portale prevede tre macro-livelli di accesso, come riassunto nella tabella seguente:
	
	\begin{table}[h]
		\centering
		\begin{tabular}{>{\bfseries}l p{10cm}}
			\toprule
			Ruolo & Permessi \\
			\midrule
			\roleadmin & Accesso completo a \textbf{tutti} i cori e a tutte le funzioni (supervisione totale). \\
			\roledirettore / \roleresponsabile & Gestione completa del \textbf{proprio coro} di appartenenza: registro presenze, gestione utenti del coro, gestione canti, bacheca. Le funzioni sono identiche per entrambi i ruoli. \\
			\rolecorista / \rolestrumentista & Area personale: profilo, visione delle proprie presenze, bacheca e calendario. Le funzioni sono identiche per entrambi i ruoli. \\
			\bottomrule
		\end{tabular}
		\caption{Ruoli e permessi del portale}
		\label{tab:ruoli}
	\end{table}
	
	Le pagine amministrative (\texttt{utenti.html}, \texttt{registro.html}, \texttt{canti-admin.html}) richiedono uno dei ruoli \roleadmin, \roledirettore o \roleresponsabile.
	
	\section{Area pubblica (senza login)}
	
	\subsection{Home page di selezione coro (\texttt{/index.html})}
	
	La home page principale permette di scegliere il coro di interesse. Ogni coro ha una propria pagina dedicata con una palette colori e un Google Sheet specifico.
	
	\subsection{Repertorio canti (\texttt{/canti.html})}
	
	Il repertorio è unificato per tutti i cori e include:
	
	\begin{itemize}
		\item \textbf{Ricerca full-text} su titoli, testi e categorie.
		\item \textbf{Bonus di punteggio} per corrispondenze esatte, per dare priorità ai titoli completi.
		\item \textbf{Filtraggio per categoria liturgica}.
		\item \textbf{Modalità presentazione} con scrolling automatico, utile per la proiezione durante le prove e le celebrazioni.
	\end{itemize}
	
	\subsection{Richiesta di un nuovo canto (\texttt{/portale/aggiungi-canto.html})}
	
	Qualsiasi utente (anche non loggato) può proporre l'aggiunta di un nuovo canto al repertorio.
	
	\textbf{Campi obbligatori:}
	\begin{itemize}
		\item Nome del canto
		\item Testo
	\end{itemize}
	
	\textbf{Campi facoltativi:}
	\begin{itemize}
		\item Autore
		\item Categoria
		\item Link YouTube
		\item Numero sul libretto
		\item Note
	\end{itemize}
	
	Dopo l'invio, il richiedente viene reindirizzato a una pagina di conferma (\texttt{/system/grazie.html}).  
	L'inserimento non è automatico: ogni richiesta viene verificata manualmente dai responsabili prima di essere pubblicata.
	
	\section{Area riservata (\texttt{/portale})}
	
	Dopo il login, l'utente accede alla \textbf{dashboard} (\texttt{/portale/index.html}), dove vengono mostrate le card dinamiche in base al ruolo.
	
	\subsection{Profilo personale (\texttt{/portale/profilo.html})}
	
	Tutti gli utenti possono:
	\begin{itemize}
		\item Visualizzare e modificare i propri dati personali
		\item Indicare il proprio \textbf{ruolo vocale/strumento}
		\item Cambiare la password
	\end{itemize}
	
	\subsection{Bacheca avvisi (\texttt{/portale/bacheca.html})}
	
	La bacheca mostra gli avvisi, filtrati per:
	\begin{itemize}
		\item \textbf{Generali} (visibili a tutti)
		\item \textbf{Specifici per coro} (visibili solo ai membri di quel coro)
	\end{itemize}
	
	\subsection{Calendario (\texttt{/portale/calendario.html})}
	
	Il calendario integrato con Google Calendar mostra:
	\begin{itemize}
		\item Prove
		\item Celebrazioni
		\item Eventi specifici per coro
	\end{itemize}
	
	\subsection{Presenze (\texttt{/portale/presenze.html})}
	
	Ogni utente può visualizzare:
	\begin{itemize}
		\item Le proprie statistiche di presenza
		\item Eventualmente, per i ruoli abilitati, le statistiche dell'intero coro
	\end{itemize}
	
	\section{Funzionalità per amministratori, direttori e responsabili}
	
	Le seguenti sezioni sono visibili e utilizzabili solo dagli utenti con ruolo \roleadmin, \roledirettore o \roleresponsabile.  
	\textit{Nota: per \roledirettore{} e \roleresponsabile{} le funzioni sono assolutamente identiche.}
	
	\subsection{Registro presenze (\texttt{/portale/registro.html})}
	
	Permette di registrare le presenze alle prove e alle celebrazioni.
	
	\subsection{Gestione utenti (\texttt{/portale/utenti.html})}
	
	Consente di:
	\begin{itemize}
		\item Creare nuovi account
		\item Attivare/disattivare utenti
		\item Assegnare ruoli
	\end{itemize}
	
	\subsection{Gestione canti (\texttt{/portale/canti-admin.html})}
	
	Da qui è possibile:
	\begin{itemize}
		\item Gestire i \textbf{canti della settimana}
		\item Amministrare il \textbf{repertorio completo}
	\end{itemize}
	
	\subsection{Richiesta nuovi canti (versione admin)}
	
	Anche gli amministratori possono utilizzare la pagina di richiesta nuovi canti (\texttt{/portale/aggiungi-canto.html}), ma in più possono pubblicare direttamente i canti approvati.
	
	\section{Suggerimenti per l'uso quotidiano}
	
	\begin{itemize}
		\item \textbf{Per i coristi:} utilizza la dashboard per tenere traccia delle tue presenze, controllare il calendario delle prove e leggere gli avvisi del tuo coro.
		\item \textbf{Per i direttori/responsabili:} pianifica le prove e le celebrazioni tramite il calendario, aggiorna il repertorio settimanale e tieni traccia delle presenze.
		\item \textbf{Per l'amministratore:} hai il controllo completo su tutti i cori e su tutte le funzionalità; gestisci utenti, canti e avvisi in modo centralizzato.
	\end{itemize}
	
	\section{Design e accessibilità}
	
	Il portale è \textbf{responsive} e si adatta automaticamente a:
	\begin{itemize}
		\item Desktop
		\item Tablet
		\item Smartphone
	\end{itemize}
	
	I breakpoint principali sono a 768px e 426px.  
	Sotto i 768px viene attivato un menu hamburger.  
	Ogni coro ha una \textbf{palette dinamica} tramite variabile CSS \texttt{--c-primary}, che cambia colore in base al coro selezionato (\texttt{coro-10}, \texttt{coro-1115}, \texttt{coro-estate}).
	
	\section{Note tecniche per gli utenti avanzati}
	
	\begin{itemize}
		\item Il frontend è ospitato su \textbf{GitHub Pages}.
		\item Il backend (API REST, autenticazione) è un servizio Node.js/Express su \textbf{Render}.
		\item Il database è \textbf{MongoDB Atlas}.
		\item Le chiamate API dal frontend passano attraverso i moduli in \texttt{/script/api.js} e \texttt{/script/auth-guard.js}.
	\end{itemize}
	
	\section{Contatti e supporto}
	
	Per qualsiasi problema o richiesta di assistenza, contatta i responsabili del portale o scrivi all'indirizzo email indicato nella sezione \textbf{Privacy e note legali} (\texttt{/privacy.html}).
	
	\vspace{1cm}
	\begin{center}
		\textit{Manuale redatto in data \today}
	\end{center}
	
\end{document}